
\documentclass{article}
\usepackage{colortbl}
\usepackage{makecell}
\usepackage{multirow}
\usepackage{supertabular}

\begin{document}

\newcounter{utterance}

\centering \large Interaction Transcript for game `multimodal\_referencegame', experiment `letter\_grids', episode 0 with NEW\_Reflect\_Agent.
\vspace{24pt}

{ \footnotesize  \setcounter{utterance}{1}
\setlength{\tabcolsep}{0pt}
\begin{supertabular}{c@{$\;$}|p{.15\linewidth}@{}p{.15\linewidth}p{.15\linewidth}p{.15\linewidth}p{.15\linewidth}p{.15\linewidth}}
    \# & $\;$A & \multicolumn{4}{c}{Game Master} & $\;\:$B\\
    \hline

    \theutterance \stepcounter{utterance}  
    & & \multicolumn{4}{p{0.6\linewidth}}{
        \cellcolor[rgb]{0.9,0.9,0.9}{
            \makecell[{{p{\linewidth}}}]{
                \texttt{\tiny{[P1$\langle$GM]}}
                \texttt{You are given three images, one is called target and the other two are distractors.} \\
\texttt{Your task is to generate a referring expression that best describes the target image while distinguishing it from the two other distractor images.} \\
\texttt{The first image is the target, the second image is a distractor, and the third image is a distractor. Instruction: Describe the target image. Generate the referring expression starting with the tag "Expression: " for the given target image. Omit any other text.} \\
            }
        }
    }
    & \\ \\

    \theutterance \stepcounter{utterance}  
    & \multicolumn{4}{p{0.45\linewidth}}{
        \cellcolor[rgb]{0.95,0.95,0.95}{
            \makecell[{{p{\linewidth}}}]{
                \texttt{\tiny{[P1$|$internal]}}
                \texttt{Reflections from previous episodes:} \\
\texttt{{-}} \\
\texttt{ Failed. I did not generate a referring expression as instructed; I only acknowledged the task. Next time, I must immediately output "Expression: " followed by a descriptive phrase that uniquely identifies the target image.} \\
\texttt{{-}} \\
\texttt{ Failed. I did not generate a referring expression as instructed; I only acknowledged the task. Next time, I must directly output "Expression: " followed by the descriptive text, ignoring any preamble.} \\
\texttt{{-}} \\
\texttt{ Failed. I generated a generic description ("a red apple") that didn't distinguish the target from the distractors (which also contained red apples). I didn't analyze the key *differences* between the images. Next time, I will systematically compare all three images first, identify the unique feature(s) of the target, and then generate an expression based on that contrast.} \\
            }
        }
    }
    & & \\ \\

    \theutterance \stepcounter{utterance}  
    & \multicolumn{4}{p{0.45\linewidth}}{
        \cellcolor[rgb]{0.95,0.95,0.95}{
            \makecell[{{p{\linewidth}}}]{
                \texttt{\tiny{[P1$|$internal]}}
                \texttt{{[}attempt 1 – valid{]}} \\
\texttt{Expression: a grid with a single row of X's in the middle row, while the other two rows have a mix of X's and â{-}č symbols} \\
            }
        }
    }
    & & \\ \\

    \theutterance \stepcounter{utterance}  
    & \multicolumn{4}{p{0.6\linewidth}}{
        \cellcolor[rgb]{0.9,0.9,0.9}{
            \makecell[{{p{\linewidth}}}]{
                \texttt{\tiny{[P1$\rangle$GM]}}
                \texttt{Expression: a grid with a single row of X's in the middle row, while the other two rows have a mix of X's and â{-}č symbols} \\
            }
        }
    }
    & & \\ \\

    \theutterance \stepcounter{utterance}  
    & & & \multicolumn{2}{p{0.3\linewidth}}{
        \cellcolor[rgb]{0.9,0.9,0.9}{
            \makecell[{{p{\linewidth}}}]{
                \texttt{\tiny{[GM$|$GM]}}
                \texttt{Expression: a grid with a single row of X's in the middle row, while the other two rows have a mix of X's and â{-}č symbols} \\
            }
        }
    }
    & & \\ \\

    \theutterance \stepcounter{utterance}  
    & & \multicolumn{4}{p{0.6\linewidth}}{
        \cellcolor[rgb]{0.9,0.9,0.9}{
            \makecell[{{p{\linewidth}}}]{
                \texttt{\tiny{[GM$\rangle$P2]}}
                \texttt{You are given three images.} \\
\texttt{You are also given a referring expression that describes one of the given images.} \\
\texttt{Your task is to select the image that matches the given referring expression.} \\
\texttt{Generate only the number (in text) of the image that matches the given expression by selecting first, second, or third.} \\
\\ 
\texttt{Expression: a grid with a single row of X's in the middle row, while the other two rows have a mix of X's and â{-}č symbols} \\
\texttt{Question: Which image does the expression refer to?} \\
\texttt{Start with the tag "Answer: ", followed by your selection. Omit any other text.} \\
            }
        }
    }
    & \\ \\

    \theutterance \stepcounter{utterance}  
    & & & \multicolumn{4}{p{0.45\linewidth}}{
        \cellcolor[rgb]{0.95,0.95,0.95}{
            \makecell[{{p{\linewidth}}}]{
                \texttt{\tiny{[P2$|$internal]}}
                \texttt{Reflections from previous episodes:} \\
\texttt{{-}} \\
\texttt{ Failed. I did not select an image because the expression was ambiguous without seeing the grid images. Next time, I would request the images to be provided or clarify the grid layout before attempting to choose.} \\
\texttt{{-}} \\
\texttt{ Failed. I misinterpreted the expression, focusing on the literal 'X' shape instead of the grid cell position. I selected an image where an 'X' was in the top{-}right area of the picture but not in the top{-}right *corner of the grid*.} \\
\\ 
\texttt{Next time, I will parse the spatial description ("top{-}right corner of the grid") more precisely, ignoring the object's shape unless specified. I will check the grid structure in each image first.} \\
\texttt{{-}} \\
\texttt{ Failed. I incorrectly selected the image with the 'X' in the top{-}right cell, but the expression specified the *corner of the grid*, meaning the entire 3x3 square, not just a cell. I should have verified the grid's boundaries more carefully. Next time, I will interpret "corner of the grid" as the outermost corner region, not just any top{-}right cell.} \\
            }
        }
    }
     \\ \\

    \theutterance \stepcounter{utterance}  
    & & & \multicolumn{4}{p{0.45\linewidth}}{
        \cellcolor[rgb]{0.95,0.95,0.95}{
            \makecell[{{p{\linewidth}}}]{
                \texttt{\tiny{[P2$|$internal]}}
                \texttt{{[}attempt 1 – valid{]}} \\
\texttt{Answer: second} \\
            }
        }
    }
     \\ \\

    \theutterance \stepcounter{utterance}  
    & & & \multicolumn{4}{p{0.6\linewidth}}{
        \cellcolor[rgb]{0.9,0.9,0.9}{
            \makecell[{{p{\linewidth}}}]{
                \texttt{\tiny{[GM$\langle$P2]}}
                \texttt{Answer: second} \\
            }
        }
    }
     \\ \\

    \theutterance \stepcounter{utterance}  
    & & & \multicolumn{2}{p{0.3\linewidth}}{
        \cellcolor[rgb]{0.9,0.9,0.9}{
            \makecell[{{p{\linewidth}}}]{
                \texttt{\tiny{[GM$|$GM]}}
                \texttt{second} \\
            }
        }
    }
    & & \\ \\

\end{supertabular}
}

\end{document}
